%% chapters/assignment/intro.tex
%% 序章　交通ネットワーク配分とは

\section*{序章 行動モデルとは}
\addcontentsline{toc}{section}{序章 行動モデルとは}

\subsection*{交通行動の記述}

前章では，交通ネットワークに対して需要結果がどのように分布するかを記述する交通量配分問題について解説した．
そこでの利用者行動は主に効用最大化理論に基づいて駆動されていた．
つまり各利用者が各経路の確定的ないし確率的な効用を評価し，各々が最大の効用をもたらす経路を選択するという仮定のもとに問題を記述した．

本章では焦点を経路選択行動自体に絞り，より一般的な枠組みである\textbf{行動モデル}（Behavioral Model）について解説する．
行動モデルとは概念的な表現であり，個人の行動選択の過程を記述するモデルの総称である．
交通量配分問題が意思決定過程自体よりもむしろその結果である交通量自体やその分布に着目しているのに対し，行動モデルは個々の利用者がどのようなプロセスを経て選択を行うのかを記述することに主眼を置いている．
あるいは，前者がネットワークシステム全体の挙動に関心があるのに対し，後者は個々の利用者の行動というより粒子的な挙動に関心があるとも言える．

行動モデル理論を導入するにあたり留意すべき点として，前章で扱った交通量配分は交通ネットワーク上の均衡概念を軸に整理しやすい一方で，行動モデルは主に効用最大化理論に基づいており理論的な出自が交通工学の内部にあるわけではないことが挙げられる．
さらに行動モデルの対象とする選択行動も経路選択にとどまらず，交通手段選択，目的地選択，出発時刻選択，活動種類・時間配分など多岐に渡る上，離散選択・連続選択，選択肢相関・集合生成など論点も多様である．
このように行動モデルはその理論的な出自も対象とする選択行動も多様であるため，「行動モデル理論」という括りは厳密には存在しないとも言えよう．
本章は交通行動の記述に関する理論的な基礎を提供することを目的とし，特に効用最大化理論に基づく行動モデルを中心に解説する．

行動モデル理論がカバーする交通行動の具体例とその論点としては，以下のようなものが挙げられる．
\begin{description}
  \item[目的地選択] 
  \item[経路選択] 
  ここでは特に選択肢集合の記述が問題となる．
  \item[]  
\end{description}

\subsection*{本節の構成}

本章の構成は以下の通りである．
\begin{description}
  \item[第1節 ランダム効用最大化理論（RUT）]
   多くの行動モデルの基礎となるランダム効用最大化理論を導入する．
   また最も基本的なモデルとして多項選択モデル（Multinomial Logit Model; MNL）の定式化を行い，その特性を解説する．

  \item[第2節 選択肢相関や異質性を考慮したモデル] 
  MNLモデルでは説明できない選択肢間の相関や意思決定主体間の異質性を考慮するためのモデルを解説する．

  \item[第3節 選択肢非列挙型モデル]
  離散選択モデルはあらかじめ選択肢集合が規定されていることを前提とするが，例えばOD間に無数の経路が存在する場合など，選択肢の列挙は必ずしも容易でない．
  本節では，この課題に対処するために提案されてきた選択肢非列挙型モデルを解説する．
  
  \item[第4節 アクティビティモデル]
  行動モデルが記述する選択行動は単独のトリップに限らない．
  最も包括的な行動表現は活動経路（Activity Path）であろう．
  本節では活動経路の選択行動を記述するActivity Based Models (ABMs)を解説し，付随する離散連続モデルについても触れる．
  
  \item [第5節 行動モデルの推定]
  行動モデルは，実際の交通行動データを用いてモデルパラメータを推定して初めて解釈や政策分析に利用できる．
  本節では，行動モデルの推定方法について，特に最大尤度法を中心に解説する．
\end{description}

% 例えば東京から大阪までの経路選択行動を考えてみる．
% まず利用者は交通手段の選択を飛行機，新幹線，夜行バス，自動車，自転車，徒歩などの選択肢から行う．
% 続いて各交通手段ごとに複数の経路が存在する場合には，さらに経路選択を行う．
% ここでは高々2つの選択行動として述べたが，実際の行動選択はより複雑である．

% 第一に，選択行動の構造自体の問題がある．
% 直前の例では「交通手段選択」→「経路選択」という階層構造を仮定したが，選択構造は個別の問題に依存して様々である．
% 第二に，選択肢集合の記述の問題がある．
% 自動車免許を保有していない人にとっては自動車は選択肢になりえないし，一部の愛好家を除いて自転車や徒歩は選択肢になりえない．
% さらに自動車を選択した場合も，高速道路の乗り降り地点や実際のODと高速道路までをつなぐ一般道区間のアクセス・イグレス経路には無数の選択肢が存在する．
% しかしながら無限の選択肢を想定して選択行動しているとは考えにくく，この問題も個別の設定に依存して様々な仮定が置かれる．
% このように選択肢集合はその定義自体が問題に依存して様々であり，計算量の観点からも適切な定式化が求められる．

% また前章は経路経路という離散的な選択肢を前提としているが，例えば1日の中の活動パタンの決定や観光地などにおける回遊行動は活動種類・活動場所等の離散選択と時間配分等の連続選択を同時に記述する必要があろう．
% このような選択行動に対しても多様な行動モデルが提案されている．


% \begin{description}
%   \item[街路・ネットワーク設計]
    
%     単純な容量増強が時に逆効果となることを示すBraess のパラドックス
%     （\sref{sec:ch3-ue} \ssref{subsec:braess}）は，モデルによる事前検討の重要性を示す典型例である\cite{Braess1968}．
%   % \item[活動経路]
%   \item[政策評価]
%    料金政策やネットワークサービスレベルに対する
%     外部不経済（混雑の外部性）を内部化するための最適な課金額を算定する．
%     Pigouvian 税の概念を交通に適用した混雑課金（\sref{sec:ch5-so}・\sref{sec:ch7-pigou}）は，
%     ロンドン，ストックホルム，シンガポールなど世界の主要都市で実施されている\cite{Pigou1920,Vickrey1969}．
%   \item[相互作用]
%     リアルタイムの交通情報や経路誘導が利用者行動に与える影響を分析し，ネットワーク全体の効率改善策を設計する．
%   \item[避難行動]
%     避難行動は避難開始時刻選択，避難目的地選択，避難経路選択などを含む多層的な選択行動として記述できる．
%     % 特に時間に応じて到達可能な避難目的地が動的に変化すること，情報とその不確実性の評価なども考慮する必要がある．
%   \item[シミュレーションとの連携]
%     ミクロシミュレーション（\sref{sec:ch8-simulation}）における個々のエージェントの行動ルールとして利用される他，アクティビティモデルはアクティビティシミュレーションの基礎をなす．
%     % 初期条件・境界条件として配分結果が利用される\cite{Daganzo2007,Geroliminis2008}．
% \end{description}

% \subsection*{行動モデルの分類}

% 交通量配分モデルは，（i）旅行時間がリンク交通量に依存するか否か（フロー依存性），
% および（ii）利用者の経路選択が確定的か確率的かの観点から，
% 大きく4種類に分類できる（図\ref{fig:model_overview}）\cite{土木学会1998,Sheffi1985}．

% \begin{description}
%   \item[All-or-Nothing 配分（AoN）]
%     交通量によらず旅行時間が一定（フロー非依存）と仮定し，
%     各利用者が最短経路にすべての交通量を集中させる手法（\sref{sec:ch2-uncongested} \ssref{subsec:aon}）．
%     計算は単純だが，混雑の非線形性を無視する．
%   \item[確率的配分]
%     フロー非依存の下で，利用者が複数の経路からロジット等の確率モデルに従って選択する手法
%     （\sref{sec:ch2-uncongested} \ssref{subsec:stochastic-assignment}・\ssref{subsec:dial-algorithm}）．経路の分散利用を反映できる．
%   \item[利用者均衡配分（User Equilibrium, UE）]
%     旅行時間がフロー依存でも，個々の利用者が自己の旅行時間を
%     最小化した結果，均衡状態が実現すると仮定する手法（\sref{sec:ch3-ue}）．
%     Wardropの第1原則を満たす配分であり，実務で最も広く用いられる\cite{Wardrop1952}．
%   \item[確率的利用者均衡配分（Stochastic User Equilibrium, SUE）]
%     フロー依存の下で，旅行時間の認知に確率的誤差が存在することを考慮した均衡配分
%     （\sref{sec:ch4-sue}）．UEより現実的な行動仮定を持つが，計算コストは増大する\cite{Daganzo1977}．
% \end{description}

% さらに，社会全体の総旅行時間を最小化するシステム最適配分（\sref{sec:ch5-so}），
% 均衡へ至る動的調整過程を記述する進化ゲーム理論的アプローチ（\sref{sec:ch6-evogame}），
% 混雑課金の経済理論（\sref{sec:ch7-pigou}），そして各種シミュレーション手法（\sref{sec:ch8-simulation}）へと
% 議論は展開していく．

% \subsection*{行動モデルの特長}

% 行動モデルが交通計画や政策評価において広く用いられてきた理由は，
% 以下の4つの特長による\cite{土木学会1998}．

% \begin{enumerate}
%   \item \textbf{論理の明快さ}：
%     「利用者は自己の旅行時間を最小化する」という単純な行動仮定から出発し，
%     直感的で説得力のある政策提言が得られる．

%   \item \textbf{理論解の一意性}：
%     適切な条件（単調増加なリンクコスト関数）の下でUE解は一意に存在し，
%     誰が計算しても同じ結果が得られる（\sref{sec:ch3-ue} \ssref{subsec:beckmann}）．
%     これは意思決定支援ツールとしての信頼性を担保する．

%   \item \textbf{ミクロ経済学との整合性}：
%     消費者余剰・生産者余剰の概念を用いた余剰分析（\sref{sec:ch5-so} \ssref{subsec:surplus-analysis}）や，
%     混雑課金の最適設計（\sref{sec:ch7-pigou}）は，ミクロ経済理論と厳密に整合しており，
%     福祉評価・費用便益分析に直接応用できる．

%   \item \textbf{複雑モデルの比較基準}：
%     エージェントベースシミュレーションや機械学習モデルなど，
%     より複雑な新しいアプローチの結果を評価する際，
%     均衡配分モデルの解は理論的な基準点（ベンチマーク）としての役割を果たす．
% \end{enumerate}


% また，本節末には補足として，黄金分割法（補足A），
% All-or-Nothing 配分の効率的アルゴリズム（補足B），
% および利用者均衡と Nash 均衡の等価性の証明（補足C）を付す\cite{Nash1950,Nash1951}．